\chapter*{ТЕОРЕТИЧЕСКАЯ ЧАСТЬ}
\addcontentsline{toc}{chapter}{ТЕОРЕТИЧЕСКАЯ ЧАСТЬ}

\section*{Вопросы из лабораторной 8}

\subsection*{1. Дайте определение электронной цифровой подписи}

\textbf{Электронная цифровая подпись (ЭЦП)} — это криптографическое преобразование набора данных, которое позволяет получателю данных убедиться в их целостности и подлинности, а также защищает от отказа отправителя от авторства.

ЭЦП обеспечивает три основные функции безопасности:

\subsubsection*{1.1. Аутентификация}
Подтверждение авторства документа. Подпись может быть создана только обладателем секретного (приватного) ключа, что гарантирует подлинность отправителя.

\subsubsection*{1.2. Целостность данных}
Любое изменение в подписанном документе делает подпись недействительной. Даже изменение одного бита в данных приводит к несовпадению вычисленного хэша с расшифрованным из подписи.

\subsubsection*{1.3. Неотказуемость (non-repudiation)}
Отправитель не может отрицать факт создания подписи, так как только он обладает приватным ключом, необходимым для её формирования.

\subsubsection*{1.4. Принцип работы ЭЦП}

Процесс создания и проверки ЭЦП включает следующие этапы:

\textbf{Создание подписи:}
\begin{enumerate}
    \item Вычисляется хэш-функция от документа: $H = hash(M)$
    \item Хэш шифруется приватным ключом отправителя: $S = E_{priv}(H)$
    \item Подпись $S$ прикрепляется к документу $M$
\end{enumerate}

\textbf{Проверка подписи:}
\begin{enumerate}
    \item Получатель вычисляет хэш документа: $H_1 = hash(M)$
    \item Расшифровывает подпись публичным ключом отправителя: $H_2 = D_{pub}(S)$
    \item Сравнивает хэши: если $H_1 = H_2$, подпись верна
\end{enumerate}

\textbf{Основные свойства:}
\begin{itemize}
    \item Подпись создаётся с использованием приватного ключа
    \item Проверка выполняется с использованием публичного ключа
    \item Невозможно создать подпись без знания приватного ключа
    \item Подпись уникальна для каждого документа
\end{itemize}

\subsection*{2. Опишите алгоритм создания цифровой подписи}

Алгоритм создания цифровой подписи состоит из нескольких этапов:

\subsubsection*{2.1. Генерация ключевой пары}

Первый этап — создание пары криптографических ключей (для асимметричного шифрования):

\textbf{Для RSA:}
\begin{enumerate}
    \item Выбираются два больших простых числа $p$ и $q$
    \item Вычисляется модуль: $n = p \times q$
    \item Вычисляется функция Эйлера: $\phi(n) = (p-1)(q-1)$
    \item Выбирается публичная экспонента $e$ (обычно 65537), взаимно простая с $\phi(n)$
    \item Вычисляется приватная экспонента $d$: $d \equiv e^{-1} \pmod{\phi(n)}$
    \item Публичный ключ: $(n, e)$
    \item Приватный ключ: $(n, d)$ (также хранятся $p, q$)
\end{enumerate}

\subsubsection*{2.2. Вычисление хэш-функции документа}

Для создания подписи необходимо вычислить хэш-функцию от исходного документа:
$$H = hash(M)$$

где $M$ — исходное сообщение (документ), $H$ — хэш-значение фиксированной длины.

\textbf{Требования к хэш-функции:}
\begin{itemize}
    \item Детерминированность (один и тот же вход даёт одинаковый хэш)
    \item Быстрое вычисление
    \item Необратимость (невозможно восстановить исходные данные по хэшу)
    \item Стойкость к коллизиям (сложно найти два разных документа с одинаковым хэшом)
    \item Лавинный эффект (малое изменение входа сильно меняет хэш)
\end{itemize}

\textbf{Примеры хэш-функций:}
\begin{itemize}
    \item SHA-256 (256-битный хэш)
    \item SHA-3 (Keccak)
    \item MD5 (устаревший, небезопасный)
    \item DJB2 (простая, быстрая, используется в данной работе для демонстрации)
\end{itemize}

\subsubsection*{2.3. Шифрование хэша приватным ключом}

Полученный хэш-значение шифруется приватным ключом владельца:
$$S = H^d \bmod n$$

где $S$ — цифровая подпись, $d$ — приватная экспонента, $n$ — модуль RSA.

Это преобразование может выполнить только владелец приватного ключа.

\subsubsection*{2.4. Прикрепление подписи к документу}

Подпись $S$ прикрепляется к исходному документу $M$. Возможны два варианта:

\textbf{Вариант 1: Подпись с сообщением}
\begin{itemize}
    \item Формат файла: $[размер\_подписи][подпись][исходные\_данные]$
    \item Размер файла увеличивается на размер подписи
    \item Этот подход используется в данной работе
\end{itemize}

\textbf{Вариант 2: Отдельная подпись}
\begin{itemize}
    \item Подпись хранится в отдельном файле
    \item Исходный файл не изменяется
    \item Требуется передача двух файлов
\end{itemize}

\subsubsection*{2.5. Алгоритм в виде псевдокода}

\begin{verbatim}
Алгоритм создания ЭЦП:

Вход: документ M, приватный ключ (n, d)
Выход: подписанный документ (M, S)

1. H = compute_hash(M)           // Хэширование документа
2. H_mod = H mod n                // Приведение к диапазону модуля
3. S = mod_pow(H_mod, d, n)       // Шифрование хэша
4. output = [size(S)][S][M]       // Формирование выходного файла
5. return output
\end{verbatim}

\textit{В данной работе для хэширования используется алгоритм DJB2, а для создания подписи — RSA с 64-битным модулем.}

\subsection*{3. Опишите алгоритм проверки цифровой подписи}

Проверка цифровой подписи — это процесс, позволяющий убедиться в подлинности и целостности документа.

\subsubsection*{3.1. Получение подписанного документа}

Получатель получает файл, содержащий:
\begin{itemize}
    \item Размер подписи (4 байта)
    \item Цифровую подпись $S$ (8 байт)
    \item Исходный документ $M$
\end{itemize}

Формат: $[sig\_size][signature][original\_data]$

\subsubsection*{3.2. Извлечение компонентов}

Из полученного файла извлекаются:
\begin{enumerate}
    \item Размер подписи для валидации формата
    \item Подпись $S$
    \item Исходные данные документа $M$
\end{enumerate}

Проверяется корректность формата файла (минимальный размер, корректность размера подписи).

\subsubsection*{3.3. Вычисление хэша оригинальных данных}

Получатель вычисляет хэш извлечённых данных:
$$H_1 = hash(M)$$

Используется та же хэш-функция, что и при создании подписи (в данной работе — DJB2).

\subsubsection*{3.4. Расшифровка подписи публичным ключом}

Подпись расшифровывается с использованием публичного ключа отправителя:
$$H_2 = S^e \bmod n$$

где $e$ — публичная экспонента, $n$ — модуль RSA.

Это даёт хэш-значение, которое было зашифровано при создании подписи.

\subsubsection*{3.5. Сравнение хэшей}

Сравниваются два хэш-значения:
\begin{itemize}
    \item $H_1$ — вычисленный хэш полученных данных
    \item $H_2$ — расшифрованный хэш из подписи
\end{itemize}

\textbf{Если $H_1 = H_2$:}
\begin{itemize}
    \item Подпись \textbf{верна}
    \item Документ не был изменён
    \item Подпись создана владельцем приватного ключа
    \item Оригинальный файл восстанавливается и сохраняется
\end{itemize}

\textbf{Если $H_1 \neq H_2$:}
\begin{itemize}
    \item Подпись \textbf{неверна}
    \item Возможные причины:
    \begin{itemize}
        \item Документ был изменён после подписания
        \item Подпись повреждена
        \item Подпись создана другим ключом
        \item Использована неправильная хэш-функция
    \end{itemize}
    \item Файл не сохраняется (компрометирован)
\end{itemize}

\subsubsection*{3.6. Алгоритм в виде псевдокода}

\begin{verbatim}
Алгоритм проверки ЭЦП:

Вход: подписанный документ, публичный ключ (n, e)
Выход: статус проверки (valid/invalid)

1. Извлечь из файла:
   sig_size = read_uint32(file)
   S = read_uint64(file)
   M = read_remaining_data(file)

2. Проверить формат:
   if sig_size != 8 then return INVALID

3. H1 = compute_hash(M)           // Вычислить хэш данных
4. H1_mod = H1 mod n               // Привести к диапазону модуля

5. H2 = mod_pow(S, e, n)           // Расшифровать подпись

6. if H1_mod == H2 then
       save_file(M)                // Сохранить оригинал
       return VALID                // Подпись верна
   else
       return INVALID               // Подпись неверна
\end{verbatim}

\textit{В данной работе при неверной подписи файл не сохраняется для предотвращения использования скомпрометированных данных.}

\section*{Вопросы из лабораторной 9}

\subsection*{1. Дайте определение электронной цифровой подписи}

Данный вопрос был подробно рассмотрен в Блоке 1, пункт 1.

\subsection*{2. Дайте определение электронного сертификата}

\textbf{Электронный сертификат (цифровой сертификат)} — это электронный документ, который связывает публичный ключ с идентификационными данными его владельца и подтверждается цифровой подписью доверенного центра сертификации (Certification Authority, CA).

\subsubsection*{2.1. Структура электронного сертификата}

Сертификат обычно соответствует стандарту X.509 и содержит:

\textbf{Основные поля:}
\begin{itemize}
    \item \textbf{Версия} — версия стандарта X.509 (обычно v3)
    \item \textbf{Серийный номер} — уникальный идентификатор сертификата
    \item \textbf{Алгоритм подписи} — используемый алгоритм (RSA, ECDSA и др.)
    \item \textbf{Издатель (Issuer)} — имя центра сертификации, выдавшего сертификат
    \item \textbf{Срок действия} — период валидности (Not Before / Not After)
    \item \textbf{Субъект (Subject)} — данные владельца сертификата
    \begin{itemize}
        \item CN (Common Name) — доменное имя или ФИО
        \item O (Organization) — организация
        \item OU (Organizational Unit) — подразделение
        \item C (Country) — страна
        \item E (Email) — электронная почта
    \end{itemize}
    \item \textbf{Публичный ключ субъекта} — публичный ключ владельца
    \item \textbf{Цифровая подпись CA} — подпись центра сертификации
\end{itemize}

\textbf{Расширения (Extensions):}
\begin{itemize}
    \item Key Usage — назначение ключа (подпись, шифрование и т.д.)
    \item Extended Key Usage — расширенное использование (SSL-сервер, email и др.)
    \item Subject Alternative Name — альтернативные имена
    \item CRL Distribution Points — точки распространения списков отзыва
\end{itemize}

\subsubsection*{2.2. Назначение электронного сертификата}

Сертификат обеспечивает следующие функции:

\textbf{1. Аутентификация владельца:}
\begin{itemize}
    \item Подтверждает, что публичный ключ принадлежит конкретному лицу или организации
    \item Предотв��ащает атаки типа <<человек посередине>> (MITM)
\end{itemize}

\textbf{2. Доверие к публичному ключу:}
\begin{itemize}
    \item Центр сертификации проверяет владельца перед выдачей сертификата
    \item Подпись CA гарантирует подлинность связи ключа с владельцем
\end{itemize}

\textbf{3. Управление жизненным циклом ключей:}
\begin{itemize}
    \item Контроль срока действия ключей
    \item Возможность отзыва скомпрометированных сертификатов
\end{itemize}

\subsubsection*{2.3. Виды сертификатов}

\textbf{По назначению:}
\begin{itemize}
    \item SSL/TLS сертификаты — для защиты веб-сайтов (HTTPS)
    \item Code Signing — для подписи программного обеспечения
    \item Email сертификаты (S/MIME) — для защиты электронной почты
    \item Client сертификаты — для аутентификации пользователей
\end{itemize}

\textbf{По уровню проверки (для SSL):}
\begin{itemize}
    \item DV (Domain Validation) — проверяется только владение доменом
    \item OV (Organization Validation) — проверяется организация
    \item EV (Extended Validation) — расширенная проверка, зелёная строка в браузере
\end{itemize}

\subsubsection*{2.4. Цепочка доверия}

Сертификаты образуют иерархическую цепочку доверия:

\begin{verbatim}
Корневой CA (Root CA)
    └─ Промежуточный CA (Intermediate CA)
        └─ Конечный сертификат (End-Entity Certificate)
\end{verbatim}

\textbf{Принцип работы:}
\begin{itemize}
    \item Корневой CA самоподписанный, встроен в операционную систему или браузер
    \item Промежуточный CA подписан корневым CA
    \item Конечный сертификат подписан промежуточным CA
    \item При проверке строится цепочка до доверенного корневого CA
\end{itemize}

\textbf{Отзыв сертификатов:}
\begin{itemize}
    \item CRL (Certificate Revocation List) — список отозванных сертификатов
    \item OCSP (Online Certificate Status Protocol) — онлайн-проверка статуса
\end{itemize}

\textit{В данной работе не реализуется полноценная инфраструктура PKI и выдача сертификатов, однако используются ключевые принципы асимметричной криптографии: генерация пары ключей (публичный/приватный), подпись приватным ключом и проверка публичным ключом.}

\subsection*{3. Какие существуют модели организации инфраструктуры электронных сертификатов?}

\textbf{Инфраструктура открытых ключей (PKI — Public Key Infrastructure)} — это комплекс программно-аппаратных средств, организационных мер и процедур для управления жизненным циклом электронных сертификатов и ключей.

Существует несколько моделей организации PKI:

\subsubsection*{3.1. Иерархическая модель (Hierarchical PKI)}

\textbf{Структура:}
\begin{verbatim}
         Корневой CA (Root CA)
              |
    +---------+---------+
    |         |         |
  CA-1      CA-2      CA-3    (Промежуточные CA)
    |         |         |
  cert1    cert2     cert3    (Конечные сертификаты)
\end{verbatim}

\textbf{Характеристики:}
\begin{itemize}
    \item Древовидная иерархия центров сертификации
    \item Корневой CA — единственный источник доверия
    \item Промежуточные CA делегируют полномочия вниз по иерархии
    \item Строгая централизация управления
\end{itemize}

\textbf{Преимущества:}
\begin{itemize}
    \item Простота управления и администрирования
    \item Ясная цепочка доверия
    \item Легко масштабируется вертикально
    \item Высокая безопасность корневого CA (может быть offline)
\end{itemize}

\textbf{Недостатки:}
\begin{itemize}
    \item Единая точка отказа (компрометация Root CA критична)
    \item Сложность горизонтального масштабирования
    \item Жёсткая структура подчинённости
\end{itemize}

\textbf{Применение:}
\begin{itemize}
    \item Корпоративные PKI
    \item Государственные инфраструктуры
    \item Военные и закрытые системы
\end{itemize}

\subsubsection*{3.2. Сетевая модель (Mesh/Network PKI)}

\textbf{Структура:}
\begin{verbatim}
  CA-1 <──> CA-2
   ↕         ↕
  CA-3 <──> CA-4

Все CA связаны между собой перекрёстными сертификатами
\end{verbatim}

\textbf{Характеристики:}
\begin{itemize}
    \item Все CA находятся на одном уровне (peer-to-peer)
    \item Отсутствует единый корневой центр
    \item CA взаимно сертифицируют друг друга (кросс-сертификация)
    \item Децентрализованная модель доверия
\end{itemize}

\textbf{Преимущества:}
\begin{itemize}
    \item Высокая отказоустойчивость (нет единой точки отказа)
    \item Гибкость в организации доверия
    \item Возможность сотрудничества независимых организаций
\end{itemize}

\textbf{Недостатки:}
\begin{itemize}
    \item Сложность управления и администрирования
    \item Необходимость поддерживать множество доверительных отношений
    \item Сложность построения и проверки цепочек доверия
    \item Масштабируемость: $O(n^2)$ связей для $n$ центров
\end{itemize}

\textbf{Применение:}
\begin{itemize}
    \item Федеративные системы (несколько независимых организаций)
    \item Межкорпоративное взаимодействие
    \item Системы B2B (business-to-business)
\end{itemize}

\subsubsection*{3.3. Гибридная модель (Hybrid PKI)}

\textbf{Структура:}
\begin{verbatim}
      Root CA-1        Root CA-2
         |                 |
    +----+----+       +----+----+
    |         |       |         |
  CA-A     CA-B <──> CA-C     CA-D
    |         |       |         |
  cert     cert     cert      cert

Иерархия внутри доменов + кросс-сертификация между доменами
\end{verbatim}

\textbf{Характеристики:}
\begin{itemize}
    \item Комбинация иерархической и сетевой моделей
    \item Иерархия внутри каждой организации
    \item Кросс-сертификация между организациями
    \item Гибкость в настройке доверительных отношений
\end{itemize}

\textbf{Преимущества:}
\begin{itemize}
    \item Баланс между простотой управления и отказоустойчивостью
    \item Автономность внутри организаций
    \item Возможность межорганизационного взаимодействия
    \item Оптимальна для крупных распределённых систем
\end{itemize}

\textbf{Недостатки:}
\begin{itemize}
    \item Повышенная сложность по сравнению с чисто иерархической моделью
    \item Требует координации политик безопасности между организациями
\end{itemize}

\textbf{Применение:}
\begin{itemize}
    \item Крупные корпорации с филиалами
    \item Государственные системы с региональными подразделениями
    \item Отраслевые решения (банковский сектор, здравоохранение)
\end{itemize}

\subsubsection*{3.4. Модель моста (Bridge CA Model)}

\textbf{Структура:}
\begin{verbatim}
  CA-1       CA-2       CA-3
    |          |          |
    +--- Bridge CA -------+

Bridge CA связывает независимые иерархии
\end{verbatim}

\textbf{Характеристики:}
\begin{itemize}
    \item Специальный промежуточный CA (Bridge CA) соединяет независимые PKI
    \item Каждая организация сохраняет свою иерархию
    \item Bridge CA выполняет роль <<посредника>> доверия
\end{itemize}

\textbf{Преимущества:}
\begin{itemize}
    \item Упрощение кросс-сертификации (вместо $O(n^2)$ связей — $O(n)$)
    \item Автономность участников
    \item Централизованная точка для межорганизационного доверия
\end{itemize}

\textbf{Недостатки:}
\begin{itemize}
    \item Bridge CA становится критичной точкой инфраструктуры
    \item Требует согласования политик между всеми участниками
\end{itemize}

\textbf{Применение:}
\begin{itemize}
    \item Федеральные PKI (например, Federal Bridge CA в США)
    \item Объединение корпоративных PKI в рамках альянсов
\end{itemize}

\subsubsection*{3.5. Сравнительная таблица моделей}

\begin{center}
\begin{tabular}{|l|c|c|c|c|}
\hline
\textbf{Критерий} & \textbf{Иерархия} & \textbf{Сеть} & \textbf{Гибрид} & \textbf{Мост} \\
\hline
Управление & Простое & Сложное & Средне & Средне \\
\hline
Отказоустойчивость & Низкая & Высокая & Средняя & Средняя \\
\hline
Масштабируемость & Вертикаль & Ограничена & Высокая & Высокая \\
\hline
Сложность цепочек & Низкая & Высокая & Средняя & Средняя \\
\hline
Централизация & Да & Нет & Частично & Частично \\
\hline
\end{tabular}
\end{center}

\textit{В данной работе реализованы базовые принципы асимметричной криптографии (генерация ключей, подпись, проверка), что является основой для построения любой из описанных моделей PKI.}

